% TOP_PROBLEM: {"id": 30006870, "title": "Shelah's Conjecture on NIP Fields", "category_id": 18, "category": {"id": 18, "name": "logic", "display_name": "Logic", "description": "Problems in mathematical logic, model theory, proof theory, and finite model theory.", "slug": "logic", "order_index": 18, "created_at": "2026-07-31T15:26:25.671Z"}, "status": "open"}
% FIRST_POSTED: 2026-09-27
% !TeX program = lualatex
% ATTEMPT_STATUS: unresolved
\documentclass[11pt]{article}
\usepackage[margin=25mm]{geometry}
\usepackage{fontspec}
\setmainfont{Latin Modern Roman}
\usepackage{amsmath,amssymb,mathtools}
\usepackage{unicode-math}
\setmathfont{Latin Modern Math}
\usepackage{xurl,hyperref}
\hypersetup{colorlinks=true,urlcolor=blue}
\setcounter{secnumdepth}{0}
\setlength{\parindent}{0pt}
\setlength{\parskip}{5pt}
\setlength{\emergencystretch}{3em}
\begin{document}
\section*{\texorpdfstring{Shelah's Conjecture on NIP Fields}{Shelah's Conjecture on NIP Fields}}\label{30006870}
\subsection{Definitions and mathematical statement}
A formula $\varphi(x;y)$ has the independence property if, in models of the theory, it can realize every subset of arbitrarily large finite sets: there are $a_1,\ldots,a_n$ and parameters $b_S$ with $\varphi(a_i;b_S)$ exactly when $i\in S$. A theory is NIP if no formula has this property. For an infinite field $K$ in the ring language, the selected conjecture asserts that NIP implies one of: $K$ is separably closed; $K$ is real closed; or $K$ admits a nontrivial henselian valuation. A valuation is henselian when its valuation ring satisfies Hensel's lifting of simple residue-field roots; nontrivial means its valuation ring is not all of $K$. No definability of the valuation is demanded by this formulation. The all-NIP-field quantifier is not replaced by finite-rank or otherwise restricted subclasses.
\par\smallskip\textit{Formulation and concept references:} \href{https://arxiv.org/abs/1911.00309}{[S1]}.

\subsection{Short English statement}
Must every infinite NIP field be separably closed, real closed, or carry a nontrivial henselian valuation?
\subsection{Research attempt: separating NIP from the missing valuation construction}
\textbf{Status and variants.} The statement concerns every infinite NIP field in the pure language of rings. It does not assume finite dp-rank, and it asks for a nontrivial henselian valuation without requiring that valuation to be definable. Johnson's primary paper [E1] proves the dp-finite case, with a definability conclusion in the relevant nonclosed cases. The characterization in [S1] concerns fields already equipped with a henselian valuation; its discussion of all NIP fields is conditional on the conjectural existence statement. These results were checked on 27 September 2026 and are not promoted here to a proof for arbitrary NIP fields.

The attempt develops three elementary parts of a possible strategy: permanence of NIP under field constructions that really are interpretations, the algebraic data needed to obtain a valuation, and the extra argument needed to make that valuation henselian. The gap between these parts is substantial.

\textbf{Lemma 341.1 (finite extensions preserve NIP).} If a field $K$ is NIP and $L/K$ is a finite field extension, then the pure field $L$ is NIP.

\emph{Proof.} Choose a $K$-basis $e_1,\ldots,e_d$ of $L$. Identify its underlying set with $K^d$. Addition is coordinatewise, while multiplication is given by finitely many structure constants
\[
 e_i e_j=\sum_{r=1}^d c_{ijr}e_r,
 \qquad c_{ijr}\in K.
\]
Thus every ring formula over $L$ translates to a formula over $K$ in tuples of $d$ coordinates, using these finitely many constants as parameters. Any pattern witnessing the independence property in $L$ would translate to the same pattern in $K$. Naming finitely many parameters does not destroy NIP: a formula with those parameters is an instance of a formula with additional parameter variables, and an independence pattern for the instance is one for the original formula. This proves the claim. $\square$

The same translation argument applies to definable subfields, allowing parameters, and to reducts. The direction for reducts is important: forgetting a valuation predicate cannot introduce new formulas, but adding such a predicate can introduce them. The NIP of the pure field cannot simply be transferred to an arbitrarily chosen valued-field expansion.

\textbf{An algebraic target for the construction.} A subring $\mathcal O\subseteq K$ containing $1$ is a valuation ring if
\[
 \forall x\in K^\times\quad x\in\mathcal O\ \text{or}\ x^{-1}\in\mathcal O.
 \tag{341.1}
\]
Here is a direct reconstruction, useful for checking what a proposed ring of ``small elements'' actually establishes. Let
\[
 \mathfrak m=\{a\in\mathcal O:a=0\text{ or }a^{-1}\notin\mathcal O\}.
\]
These are the nonunits of $\mathcal O$. If $a,b\in\mathfrak m$ are nonzero, (341.1) applied to $a/b$ shows that one of $a,b$ is an $\mathcal O$-multiple of the other. Suppose $b=ac$. Then $a+b=a(1+c)$ cannot be a unit, since a product can be a unit only if both factors are units. Also multiplication of a nonunit by any element of $\mathcal O$ gives a nonunit or zero. Hence $\mathfrak m$ is an ideal. Every element outside it is a unit, making it the unique maximal ideal.

The quotient group $\Gamma_v=K^\times/\mathcal O^\times$ becomes an ordered abelian group by declaring
\[
 v(x)\geq v(y)\quad\Longleftrightarrow\quad x/y\in\mathcal O.
\]
This is well defined under multiplication by units; (341.1) gives totality, and membership in $\mathcal O$ gives transitivity. If $v(x)\geq v(y)$, then $x+y=y(1+x/y)$, so
\[
 v(x+y)\geq\min\{v(x),v(y)\}
\]
when $x+y\ne0$. With $v(0)=\infty$, this is a valuation. It is nontrivial exactly when $\mathcal O\ne K$.

Thus a construction must provide much more than an additive subgroup or a proper subring. It must establish (341.1), properness, and eventually a lifting property for roots. Intersections of candidate rings need not preserve the total comparison in (341.1), so intersecting all vaguely compatible neighborhoods is not automatically a valuation construction.

\textbf{Lemma 341.2 (a complete discrete case of henselianity).} Let $K$ be complete for a discrete valuation $v:K^\times\to\mathbb Z$, with valuation ring $\mathcal O$ and maximal ideal $\mathfrak m$. If $f\in\mathcal O[X]$ and $a_0\in\mathcal O$ satisfy
\[
 f(a_0)\in\mathfrak m,\qquad f'(a_0)\in\mathcal O^\times,
\]
then there is a unique $a\in a_0+\mathfrak m$ with $f(a)=0$.

\emph{Proof.} Unless a root has already been found, define
\[
 a_{n+1}=a_n-\frac{f(a_n)}{f'(a_n)}.
\]
For $a,h\in\mathcal O$, polynomial expansion gives
\[
 f(a+h)=f(a)+hf'(a)+h^2R(a,h),\qquad
 f'(a+h)-f'(a)\in h\mathcal O,
\]
where $R(a,h)\in\mathcal O$. If $r_n=v(f(a_n))\geq1$ and $f'(a_n)$ is a unit, the correction $h_n$ has valuation $r_n$. The two identities show that $f'(a_{n+1})$ remains a unit and $r_{n+1}\geq2r_n$, unless the new value is zero. Hence $r_n\geq2^n$, and the corrections tend to zero rapidly. Completeness gives a limit $a\in a_0+\mathfrak m$, and continuity gives $f(a)=0$.

For uniqueness, if $a,b\in a_0+\mathfrak m$, factor
\[
 f(a)-f(b)=(a-b)U,
\]
where the polynomial difference quotient $U$ is congruent to $f'(a_0)$ modulo $\mathfrak m$. It is therefore a unit. Two roots in this residue class must be equal. $\square$

This proves the familiar simple-root form of Hensel's lemma for complete discretely valued fields. It establishes the required kind of root lifting once those stronger topological assumptions are available; it does not construct such a topology on an arbitrary NIP field.

\textbf{Two stress tests on the proposed route.} First, having a nontrivial valuation is not enough. On $\mathbb Q$ with its $3$-adic valuation, the polynomial $f(X)=X^2-7$ has a simple root $1$ modulo $3$, since $f(1)=-6$ and $f'(1)=2$. It has no rational root. Thus this valuation on $\mathbb Q$ is not henselian. Passing to a completion would produce a different field; it would not establish henselianity of the original field in the statement.

Second, henselianity does not by itself imply NIP. The complete discretely valued field $\mathbb F_p((t))$ is henselian by Lemma 341.2, but is not NIP. Here we use the primary theorem of Kaplan--Scanlon--Wagner [E2]: an infinite NIP field of characteristic $p$ has no nontrivial Artin--Schreier extension, equivalently the map $x\mapsto x^p-x$ is surjective.

For this example the failure of surjectivity is elementary. The equation
\[
 x^p-x=t^{-1}
 \tag{341.2}
\]
has no solution in $\mathbb F_p((t))$. If $v(x)\geq0$, its left side has nonnegative valuation. If $v(x)<0$, then $pv(x)<v(x)$, so the two terms have unequal valuations and
\[
 v(x^p-x)=pv(x),
\]
which is divisible by $p$ and cannot equal $-1$. The polynomial $X^p-X-t^{-1}$ therefore has no root. Its roots in a splitting field differ by elements of $\mathbb F_p$, and its Galois group is a subgroup of the additive group of order $p$; consequently absence of a root makes it an irreducible Artin--Schreier polynomial of degree $p$. This contradicts the necessary NIP condition from [E2]. The external theorem is used only for that necessary condition, not as a construction of a valuation.

\textbf{Why finite rank is an additional hypothesis.} NIP bounds the shattering complexity of each fixed formula. It does not, from its definition alone, give a single finite numerical rank controlling every family of definable sets. An argument that chooses a maximal finite collection of independent definable directions must justify the finiteness of that collection. In the dp-finite theorem [E1], a finite rank hypothesis supplies precisely additional structure that is absent from the selected statement. Dropping that hypothesis from the beginning of a proof would silently replace a theorem about one subclass by the open general assertion.

Similarly, adjoining parameters does not create a valuation for free. Lemma 341.1 says a specified interpretation preserves NIP; it does not say every abstract subring of an NIP field is definable, or that a definable topology must have a henselian valuation ring. Each of these would require a separate theorem with its own hypotheses.

\textbf{Checks and outcome.} The accompanying exact diagnostic implements Newton lifting of $X^2-7$ from $1$ modulo $3$ to increasing powers of $3$, and checks the valuation obstruction in (341.2) for several primes and negative valuations. Finite modular lifts illustrate the complete-field argument; they do not turn the resulting $3$-adic root into a rational root. No finite experiment is presented as a test for NIP of an infinite field.

The proved components establish useful permanence, reconstruct a valuation from a total ring comparison, and prove henselianity under completeness. The first unresolved step is constructing a proper ring satisfying (341.1), with enough control to obtain henselianity, for an arbitrary infinite NIP field outside the separably closed and real closed alternatives. The dp-finite theorem and the characterization of already henselian NIP fields do not remove that step. This attempt leaves the selected conjecture unresolved.

\subsection{Sources}
\begin{itemize}
\item[S1]Characterizing NIP henselian fields (arXiv abstract). Accessed 2026-09-01.
\url{https://arxiv.org/abs/1911.00309}.
\textit{Formulation source.}
\item[E1]Will Johnson, \emph{Dp-finite fields VI: the dp-finite Shelah conjecture}, arXiv:2005.13989 (2020).
\url{https://arxiv.org/abs/2005.13989}.
\textit{Primary source consulted 27 September 2026; the finite dp-rank hypothesis is essential to the result cited here.}
\item[E2]Itay Kaplan, Thomas Scanlon, and Frank O. Wagner, \emph{Artin--Schreier extensions in NIP and simple fields}, Israel Journal of Mathematics 185 (2011), 141--153; arXiv:1009.5421.
\url{https://arxiv.org/abs/1009.5421}.
\textit{Primary source consulted 27 September 2026; used for Artin--Schreier closedness of infinite NIP fields in positive characteristic.}
\end{itemize}
\end{document}
